\chapter{Methodology}

\section{The setup}

A complex scalar field of mass $m$ and charge $e$ obeys the Klein-Gordon equation
\begin{align}
	\left( D_\mu D^\mu + m^2 \right)\phi = 0.
	\label{eq:TDKGE}
\end{align}
Confined to an interval of size normalized to 1, we consider Robin boundary conditions i.e.~boundary conditions of the form of 
\begin{align}
	n_\mu D^{\mu} \phi+ \xi(\sigma)	 \phi = 0
\end{align}
with $\sigma$ the boundary of the region considered.
We are however, not interested in (anti) periodic boundary conditions, as these would translate into a `particle' that gets accelerated every time it loops around the interval.

Under the Coulomb gauge, the electromagnetic potential corresponding to a constant electric field is
\begin{align}
	A_0 = -\lambda \left(z - \frac{1}{2}\right) \,\,\,\, A_1 = 0.
\end{align}

Assuming the separation ansatz 
\begin{align}
	\phi(t, z) = \phi_n(z) e^{-i\omega_n t}
\end{align}
equation \eqref{eq:TDKGE} can be rewritten as the time independent Klein-Gordon equation
\begin{align}
	\partial_1^2 \phi_n &=
	\left( - (\omega_n - e A_0)^2 + m^2 \right) \phi_n\\
			    \label{eq:TIKGE}
			    &=\left( - \left(\omega_n + \lambda \left( z - \frac{1}{2} \right) \right)^2 + m^2 \right) \phi_n.
\end{align}
One may notice that this equation is of the form of the Weber differential equation
\begin{align}
	\partial_1^2 \phi_n = - \left( \nu + \frac{1}{2}-\frac{1}{4}\tilde{z}^2 \right) \phi_n,
\end{align}
and can thus be solved by the parabolyc cylinder functions $D_\nu\left( \tilde{z} \right) $. \note{Add this to the appendix?}
We then write the general solution to equation \ref{eq:TIKGE} as follows
\begin{align}
	\phi_n(z) = 
	a_n D_{i \frac{m^2}{2\lambda} - \frac{1}{2}} \left( 
		\frac{1+i}{\sqrt{\lambda} }
		\left( \omega_n + \lambda \left( z-\frac{1}{2} \right)
		\right) 
	\right) 
	+
	b_n D_{-i \frac{m^2}{2\lambda} - \frac{1}{2}} \left( 
		\frac{i - 1}{\sqrt{\lambda} }
		\left( \omega_n + \lambda \left( z-\frac{1}{2} \right)
		\right) 
	\right). 
\end{align}

\note{Also talk about the 0 field solution that indeed lead to vanishing vacuum polarization.}
We have thus far considered what in \cite{Ambjorn1983} is referred to as the `external field approximation'. This implies that one does not consider the fact that the vacuum polarization induces itself another electric field. This approximation fails \cite{Ambjorn1983} in the calculation of the energy of the modes, as it yields complex values for some external field values.

\note{\cite{Schlemmer2015}: This seems to be the reason why in
	\cite{Ambjorn1983}, where the mode sum formula is employed, it is claimed that the vacuum
polarization vanishes for massless fermions in 1 + 1 spacetime dimension. As
we will see, this is not correct.}

We therefore use the ``back reaction'' approach, in which we solve the problem iteratively. Following the convention used in \cite{Ambjorn1983}, we label by $\kappa$ each of the iterations. At each step $\kappa$ we solve numerically the following set of integro-differential equations
\begin{align}
	\partial_1^2 \phi^{\kappa+1}_n &=
	\left( - (\omega_n - e A^{\kappa}_0(z))^2 + m^2 \right) \phi^{\kappa+1}_n\\
	\partial_1^2\tilde{A}^{\kappa} &= -\rho^{\kappa}
	\label{eq:poisson}
,
\end{align}
with $$A_0^{\kappa}(z) = -\lambda \left( z - \frac{1}{2} \right) + \tilde{A_0}^{\kappa}(z),$$  and $\rho^{\kappa}(z)$ the associated charge density at each step.
Since we expect to see no screening at the boundaries, \eqref{eq:poisson} has Neumann boundary conditions.

\note{\cite{Schlemmer2015}:  The supplementary electromagnetic field produced by (36) via back-
reaction is 
\begin{align}
	E_{br} = \frac{1}{2\pi}e^2 E \left[ z^2 - \left( \frac{L}{2} \right) ^2 \right] 
\end{align}
For the model to be physically viable, this should be small compared to the external electric field, or equivalently
\begin{align}
	eL \ll 1
\end{align}
}

This method removes the appearance of complex eigenvalues, as the screening effects of the vacuum should be taken into account.

One should take care when calculating the charge density, since 
\begin{align}
	\rho(z) = \vacuum{\hat{\rho}(z)} = ie\vacuum{\phi^* D_0 \phi - \phi \left( D_0 \phi \right) ^*}
\end{align}
is not well defined as it involves the product of two fields. We take care of this in the following section.

\section{Point-split renormalization}

Due to the linearity of \eqref{eq:TDKGE}, we write the quantum field as 
\begin{align}
	\phi(t, z) = 
	\sum_{n>0}^{} a_n \phi_n(z) e^{-i\omega_n t}
	+ \sum_{n<0}^{} b_n^\dagger \phi_n(z) e^{-i\omega_n t}
\end{align}
with the operators $a_n$, $b_n$ fulfilling 
\begin{align}
	\left[ a_n, a_m^\dagger \right]  = \delta_{nm} \,\,\,\,\,\,\,\,\,\,
	\left[ b_n, b_m^\dagger \right]  = \delta_{nm}
\end{align}
and all other commutators vanishing.

We define the vacuum state $\ket{0}$ as that fulfilling 
\begin{align}
	a_n\ket{0} = 0, \,b_n\ket{0} = 0.
\end{align}

To calculate the charge density $\rho(z)$ of the field, one would naïvely calculate 
\begin{align}
	\rho(z) = ie\vacuum{\phi^* D_0 \phi - \phi \left( D_0 \phi \right) ^*},
\end{align}
which is known to be ill-defined as it involves the product of two distributions.

As any physically reasonable state is assumed to be Hadamard,  \note{Need to read on this}
\begin{align}
	\rho^{\kappa}(z) = 
	\lim_{\tau \to 0}
	\sum_{n=-\infty}^{\infty} (\omega_n - A_0^{\kappa}(z)) \|\phi_n(z)\|^2 e^{-i \omega_n (\tau + i\varepsilon)} 
	+ \frac{e}{\pi^2} A_0^{\kappa}(z)
\end{align}


\section{}


